% Definition file for row def_3_14 -- "MarginalizationAK".
% Auto-generated stub by scaffold/solve_chapter.py at row start. The
% manager agent (or a worker dispatched by it) fills the body below with the
% Lean-faithful definition statement(s). Stays close to the lecture notes.
%
% This file is a sibling of `main.tex` inside `<subsection>/tex/`. The
% `subfiles` package lets it compile both standalone (resolving its
% preamble via `main.tex`) and as part of `main.tex` via `\subfile{...}`.

\documentclass[main]{subfiles}

\begin{document}
% ref: def_3_14
% title: MarginalizationAK
\def\rowref{def\_3\_14}
\phantomsection\label{def_3_14}

\begin{Def}[Marginalization a.k.a.\ latent projection on CDMGs]\label{def:G_marginalization}
    Let $G = (J, V, E, L)$ be a CDMG (def \ref{def-cdmg}) and $W \ins V$ a subset of output nodes. The \emph{marginalization of $G$ w.r.t.\ $W$}, equivalently the \emph{latent projection of $G$ onto $J \cup (V \sm W)$}, is the CDMG
    \[
        G^{V \sm W \,|\, J} \;:=\; G^{\sm W} \;:=\; \lp J^{\sm W},\, V^{\sm W},\, E^{\sm W},\, L^{\sm W} \rp,
    \]
    whose four components are defined as follows.
    \begin{enumerate}[label=\roman*.)]
        \item \emph{Input nodes.}
            \[ J^{\sm W} \;:=\; J. \]

        \item \emph{Output nodes.}
            \[ V^{\sm W} \;:=\; V \sm W. \]

        \item \emph{Directed edges.} The directed-edge set is the set of pairs witnessed by a directed walk in $G$ whose intermediate vertices lie in $W$:
            \[
                E^{\sm W} \;:=\; \lC (\ul{v}, \ol{v}) \in \lp J \cup (V \sm W) \rp \times (V \sm W) \st \Phi_E(\ul{v}, \ol{v}) \rC,
            \]
            where for $\ul{v} \in J \cup (V \sm W)$ and $\ol{v} \in V \sm W$ the directed-walk-through-$W$ predicate $\Phi_E(\ul{v}, \ol{v})$ holds if and only if there exist an integer $n \ge 1$ and a tuple $(w_0, w_1, \ldots, w_n) \in (J \cup V)^{n+1}$ such that all of the following hold:
            \begin{enumerate}[label=(\alph*)]
                \item \emph{End-points.} $w_0 = \ul{v}$ and $w_n = \ol{v}$.
                \item \emph{Intermediates lie in $W$.} $w_i \in W$ for every $i \in \{1, \ldots, n - 1\}$ (vacuous when $n = 1$).
                \item \emph{Consecutive $E$-edges.} $(w_{i-1}, w_i) \in E$ for every $i \in \{1, \ldots, n\}$.
            \end{enumerate}
            In words: $\Phi_E(\ul{v}, \ol{v})$ asserts the existence of a directed walk in $G$ from $\ul{v}$ to $\ol{v}$, in the sense of def \ref{def:walks}, item~ii., of length $n \ge 1$, all of whose intermediate vertices $w_1, \ldots, w_{n-1}$ (if any) lie in $W$.

            \emph{No $\ul{v} \neq \ol{v}$ constraint.} The condition $\ul{v} \neq \ol{v}$ is \emph{not} imposed in the set-builder for $E^{\sm W}$. Self-cycles $(v, v) \in E^{\sm W}$ are admitted exactly under the witness conditions above, namely via a directed walk through $W$ that returns to $v$.\footnote{This may introduce self-cycles into $E^{\sm W}$.}

        \item \emph{Bidirected edges.} The bidirected-edge set is the set of distinct pairs witnessed by a bifurcation in $G$ whose intermediate vertices lie in $W$:
            \[
                L^{\sm W} \;:=\; \lC (\ul{v}, \ol{v}) \in (V \sm W) \times (V \sm W) \st \ul{v} \neq \ol{v} \,\land\, \Phi_L(\ul{v}, \ol{v}) \rC,
            \]
            where for $\ul{v}, \ol{v} \in V \sm W$ the bifurcation-through-$W$ predicate $\Phi_L(\ul{v}, \ol{v})$ holds if and only if there exist an integer $n \ge 1$, an index $k \in \{1, \ldots, n\}$, and a tuple $(w_0, w_1, \ldots, w_n) \in (J \cup V)^{n+1}$ such that the data $\bigl((w_0, w_1, \ldots, w_n),\, k\bigr)$ describes a bifurcation between $\ul{v}$ and $\ol{v}$ in $G$ in the sense of def \ref{def:walks}, item~vi., \emph{and} all intermediate vertices $w_1, \ldots, w_{n-1}$ (if any, i.e.\ when $n \ge 2$) lie in $W$. Unfolding def \ref{def:walks}, item~vi., at the present row, $\Phi_L(\ul{v}, \ol{v})$ is the conjunction of:
            \begin{enumerate}[label=(\alph*)]
                \item \emph{End-points, intermediates, and walk length.} $w_0 = \ul{v}$, $w_n = \ol{v}$, $n \ge 1$, $k \in \{1, \ldots, n\}$, and $w_i \in W$ for every $i \in \{1, \ldots, n-1\}$ (vacuous when $n = 1$).
                \item \emph{End-nodes appear exactly once.} $w_0 \notin \lC w_1, w_2, \ldots, w_n \rC$ and $w_n \notin \lC w_0, w_1, \ldots, w_{n-1} \rC$. (Under the constraints of (a) this is automatic, since $\ul{v}, \ol{v} \in V \sm W$, while $w_1, \ldots, w_{n-1} \in W$ and $\ul{v} \neq \ol{v}$ is enforced by the outer set-builder.)
                \item \emph{Left directed arm} (vacuous when $k = 1$). For every $i \in \{1, \ldots, k - 1\}$,
                    \[ (w_i, w_{i-1}) \in E. \]
                    (Equivalently, in def \ref{not-cdmg} notation, $w_{i-1} \hut w_i$ in $G$ for every such $i$.)
                \item \emph{Central hinge} at index $k$ --- either a directed-fork hinge or a bidirected hinge:
                    \[ (w_k, w_{k-1}) \in E \;\;\text{or}\;\; (w_{k-1}, w_k) \in L. \]
                    (Equivalently, $w_{k-1} \hus w_k$ in $G$, in def \ref{not-cdmg} notation.)
                \item \emph{Right directed arm} (vacuous when $k = n$). For every $i \in \{k, \ldots, n - 1\}$,
                    \[ (w_i, w_{i+1}) \in E. \]
                    (Equivalently, $w_i \tuh w_{i+1}$ in $G$ for every such $i$.)
                \item \emph{End-node arrowhead constraint} (inherited from def \ref{def:walks}, item~vi., clause~(e)). When $k = n$, only the \emph{bidirected} alternative of the hinge in (d) is admitted, namely $(w_{k-1}, w_k) \in L$; the directed alternative $(w_k, w_{k-1}) \in E$ at $k = n$ is \emph{excluded}, because it would place the only arrowhead incident to $w_n$ at $w_{n-1}$ rather than at $w_n = \ol{v}$, contradicting clause~(e) of def \ref{def:walks}, item~vi. For every $k$ with $1 \le k \le n - 1$, both alternatives in (d) are admitted (the right directed arm of (e) supplies an arrowhead at $w_n$ regardless of the hinge type, and the right arm also gives $w_k$ at least one out-edge so the directed-fork hinge becomes a genuine fork with source $w_k$). Symmetrically at $w_0$: for $k \ge 2$ the left arm supplies an arrowhead at $w_0$ via $(w_1, w_0) \in E$; for $k = 1$ both hinge alternatives are edges into $w_0$ (the directed alternative $(w_1, w_0) \in E$ puts an arrowhead at $w_0$, and the bidirected alternative $(w_0, w_1) \in L$ puts arrowheads at both ends), so neither is excluded at $k = 1$.
            \end{enumerate}

            \emph{Boundary cases authorised by the conventions $w_0 := \ul{v}$, $w_n := \ol{v}$ and $k \in \{1, \ldots, n\}$.} The conventions explicitly extend the displayed index range of the bifurcation walk to its endpoints, so that all of the following configurations qualify as bifurcation-through-$W$ witnesses for $L^{\sm W}$, provided $\ul{v} \neq \ol{v}$ and every intermediate $w_i$ (for $1 \le i \le n - 1$) lies in $W$:
            \begin{itemize}
                \item $n = 1$, $k = 1$ (both arms vacuous; $k = n$, so clause~(f) admits only the bidirected hinge): the witness reduces to the single hinge edge $(w_0, w_1) = (\ul{v}, \ol{v}) \in L$ --- i.e., a direct bidirected edge $\ul{v} \huh \ol{v}$ already present in $G$ between non-marginalised nodes, with no intermediates. The directed-hinge alternative is excluded here by clause~(f).
                \item $n = 2$, $k = 1$ (left arm vacuous): one intermediate $w_1 \in W$; the hinge at index $k = 1$ is $(w_1, \ul{v}) \in E$ (the ``Y''-fork case with $w_1$ as source) or $(\ul{v}, w_1) \in L$ (the bidirected-hinge case); the right directed arm contributes $(w_1, \ol{v}) \in E$.
                \item $n = 2$, $k = n = 2$ (right arm vacuous; $k = n$, so clause~(f) admits only the bidirected hinge): one intermediate $w_1 \in W$; the left directed arm contributes $(w_1, \ul{v}) \in E$; the hinge admits only the bidirected alternative $(w_1, \ol{v}) \in L$ (the mirror ``Y'' with bidirected edge into $\ol{v}$).
            \end{itemize}
            Larger values of $n$ and $k$ are also admitted, with the same structural skeleton: a left arm of directed edges all pointing toward $\ul{v}$, a single hinge (directed-fork or bidirected, subject to clause~(f)), a right arm of directed edges all pointing toward $\ol{v}$, and all intermediate vertices in $W$.
    \end{enumerate}

    \emph{Asymmetry between (iii) and (iv) regarding self-edges (intentional; to be preserved by any formalisation).}
    The constraint $\ul{v} \neq \ol{v}$ is imposed in clause (iv) on $L^{\sm W}$ but is \emph{not} imposed in clause (iii) on $E^{\sm W}$. Consequently:
    \begin{itemize}
        \item Directed self-cycles $v \tuh v$ may appear in $E^{\sm W}$ when witnessed by a directed walk through $W$ that returns to $v$.
        \item Bidirected self-edges $v \huh v$ are \emph{excluded} from $L^{\sm W}$ outright by the $\ul{v} \neq \ol{v}$ constraint, irrespective of any bifurcation-through-$W$ that might otherwise witness one.
    \end{itemize}
    Any formalisation of this definition shall preserve this asymmetry exactly: no $\ul{v} \neq \ol{v}$ constraint is to be imposed on $E^{\sm W}$, and no relaxation of the $\ul{v} \neq \ol{v}$ constraint is to be made on $L^{\sm W}$.
\end{Def}

\end{document}
